% ============================================================
% symbols.tex
% Symbol definitions and List of Symbols for the dissertation
% "Knot Floer Homology, Folding Automata, and the Connect-Sum
%  of Trefoils"
% ============================================================

\usepackage{enumitem}
\newlist{abbrv}{itemize}{1}
\setlist[abbrv,1]{label=,labelwidth=1.6in,align=parleft,%
                  itemsep=0.15\baselineskip,leftmargin=!}

\def\symbolentry#1#2#3{\item[#2] #3}
\def\sort#1{}

% ============================================================
% Number sets and basic notation
% Use \providecommand so we don't conflict with preamble.tex
% definitions (\N, \R, \Q, \C, \O may already be defined there).
% ============================================================
\providecommand{\Z}{\ensuremath{\mathbb{Z}}}
\providecommand{\N}{\ensuremath{\mathbb{N}}}
\providecommand{\R}{\ensuremath{\mathbb{R}}}
\providecommand{\Q}{\ensuremath{\mathbb{Q}}}
\providecommand{\C}{\ensuremath{\mathbb{C}}}
\newcommand{\F}{\ensuremath{\mathbb{F}}}                % a field
\newcommand{\Ftwo}{\ensuremath{\mathbb{F}_{2}}}         % field of two elements
\newcommand{\HH}{\ensuremath{\mathbb{H}}}               % quaternions / hyperbolic
\newcommand{\T}{\ensuremath{\mathbb{T}}}                % torus

% ============================================================
% Operators
% (\rank, \tr are already declared in preamble.tex — do not redefine)
% ============================================================
\DeclareMathOperator{\FDTC}{FDTC}                       % frac. Dehn twist coeff.
\DeclareMathOperator{\HFK}{\widehat{HFK}}               % hat knot Floer homology
\DeclareMathOperator{\HFKminus}{HFK^{-}}                % minus version
\DeclareMathOperator{\CFK}{\widehat{CFK}}               % chain complex
\DeclareMathOperator{\HF}{\widehat{HF}}                 % Heegaard Floer
\DeclareMathOperator{\Mod}{Mod}                         % mapping class group
\DeclareMathOperator{\PMod}{PMod}                       % pure MCG
\DeclareMathOperator{\Homeo}{Homeo}
\DeclareMathOperator{\Diff}{Diff}
\DeclareMathOperator{\id}{id}
\DeclareMathOperator{\Spin}{Spin^{c}}
\DeclareMathOperator{\genus}{g}                         % Seifert genus
\DeclareMathOperator{\sgn}{sgn}
\DeclareMathOperator{\Tw}{Tw}                           % Dehn twist
\DeclareMathOperator{\Stab}{Stab}
\DeclareMathOperator{\Fix}{Fix}
\DeclareMathOperator{\Aut}{Aut}
\DeclareMathOperator{\Out}{Out}
\DeclareMathOperator{\Hom}{Hom}
\DeclareMathOperator{\Ker}{Ker}
\renewcommand{\Im}{\operatorname{Im}}

% Inverse macro
\newcommand{\inv}[1]{#1^{-1}}

% Restriction
\newcommand{\restrict}[2]{#1\bigr|_{#2}}

% ============================================================
% Train-track / folding-automaton-specific notation
% (\tt is a TeX primitive for typewriter font; do NOT redefine it)
% ============================================================
\newcommand{\trtrack}{\tau}                             % a train track
\newcommand{\TT}{\mathcal{T}}                           % set of train tracks
\newcommand{\stratum}[1]{\mathcal{Q}(#1)}               % stratum of differentials

% Branches
\newcommand{\bbranch}{\beta}                            % a generic branch
\newcommand{\bbar}{\bar{b}}                             % branch b with bar
% \beth is provided by amssymb (Hebrew letter) — do not redefine

% Edge labels used in case-analysis appendices
% (single-letter macros chosen to mirror the manuscript)
\newcommand{\dbar}{\bar{d}}
\newcommand{\dbara}{\bar{d}}                            % alias for \dbar (used in appendix)
\newcommand{\dabr}{\bar{d}}                             % alias for \dbar (used in appendix)
\newcommand{\hato}{\hat{o}}                             % hatted o

% Branch images under the folding map
\newcommand{\fbd}{f_\beta(d)}
\newcommand{\fbdb}{f_\beta(\bar{d})}
\newcommand{\fbo}{f_\beta(o)}
\newcommand{\fbi}{f_\beta(i)}
\newcommand{\fbs}{f_\beta(s)}
\newcommand{\fbp}{f_\beta(p)}
\newcommand{\fbb}{f_\beta(b)}
\newcommand{\fbr}{f_\beta(r)}
\newcommand{\fbe}{f_\beta(e)}                           % image of e
\newcommand{\fbar}{f_\beta}                             % the folding map itself
\newcommand{\fdb}{\fbdb}                                % alias (appears in appendix typo)
\newcommand{\fdbd}{\fbdb}                               % alias
\newcommand{\fbfb}{\fbdb}                               % alias
\newcommand{\fds}{f_\beta(s)}                           % alias for \fbs
\newcommand{\dbp}{\bar{d}_\beta(p)}                     % d-bar component used in case analysis
\newcommand{\bet}{\beta}                                % alias

% Marker symbols used in case-analysis prose
\newcommand{\circf}{\circ}                              % \circ alias
\newcommand{\cric}{\circ}                               % \circ alias (typo in mss.)
\newcommand{\timed}{\times}                             % \times alias
\newcommand{\tmes}{\times}                              % \times alias

% Absorption / rule operators
% \Abso{a}{b} expands to \Absorbed(a,b)^\times — do NOT write ^\times after it
\DeclareMathOperator{\Absorbed}{Abso}
\DeclareMathOperator{\Rule}{Rule}
\newcommand{\Abso}[2]{\Absorbed(#1, #2)^\times}
\newcommand{\Rul}[1]{\Rule(#1)^\times}

% Composite superscript glyphs used in train-track labels
\newcommand{\circm}{\substack{\circ \\[-3pt] -}}
\newcommand{\pcirc}{\substack{+ \\[-3pt] \circ}}

% Superimposed symbols (for \tcap; kept from template in case used)
\makeatletter
\newcommand{\superimpose}[2]{%
  {\ooalign{$#1\@firstoftwo#2$\cr\hfil$#1\@secondoftwo#2$\hfil\cr}}}
\makeatother
\newcommand{\tcap}{\pitchfork}

% ============================================================
% Knot Floer homology / Alexander polynomial conventions
% ============================================================
\newcommand{\Alex}[1]{\Delta_{#1}}                      % Alexander polynomial
\newcommand{\Cycl}[1]{\Phi_{#1}}                        % cyclotomic polynomial
\newcommand{\trefoil}{T_{2,3}}                          % the trefoil
\newcommand{\cable}[2]{C_{#1,#2}}                       % (p,q)-cable operator
\newcommand{\granny}{\trefoil \# \trefoil}              % granny knot

% ============================================================
% List of Symbols
% ============================================================
\newcommand{\listofsymbols}{
\chapter*{List of Symbols}
\addcontentsline{toc}{chapter}{List of Symbols}
\begin{abbrv}

  % ---- Number systems and fields ----
  \symbolentry{Z}{$\Z$}{The integers}
  \symbolentry{Ftwo}{$\Ftwo$}{The field with two elements; coefficient ring for $\HFK$}
  \symbolentry{Q}{$\Q$}{The rational numbers}
  \symbolentry{R}{$\R$}{The real numbers}
  \symbolentry{C}{$\C$}{The complex numbers}

  % ---- Surfaces and mapping class groups ----
  \symbolentry{S}{$\Sigma_{g,n}$}{Orientable surface of genus $g$ with $n$ boundary components}
  \symbolentry{D}{$D_n$}{The $n$-punctured disk}
  \symbolentry{Mod}{$\Mod(\Sigma)$}{Mapping class group of the surface $\Sigma$}
  \symbolentry{PMod}{$\PMod(\Sigma)$}{Pure mapping class group fixing punctures setwise}
  \symbolentry{Tw}{$\Tw_\gamma$}{Right-handed Dehn twist about the simple closed curve $\gamma$}

  % ---- Train tracks and folding automata ----
  \symbolentry{tt}{$\tau$}{A train track on a surface}
  \symbolentry{ttbr}{$\beta$}{A branch of a train track}
  \symbolentry{ttsm}{$\stratum{\kappa}$}{Stratum of quadratic differentials of singularity profile $\kappa$}
  \symbolentry{TT}{$\TT$}{The set of standard train tracks in a fixed stratum}
  \symbolentry{ttf}{$f_\beta$}{The folding (train-track) map of a pseudo-Anosov $\beta$}
  \symbolentry{fbe}{$f_\beta(e)$}{Image of a branch $e$ under the folding map}
  \symbolentry{Aut}{$\Aut(\tau)$}{Folding automaton on standard train tracks}
  \symbolentry{Rule}{$\Rule(\beta)^\times$}{Folding rule applied to $\beta$ on the disk}
  \symbolentry{Abso}{$\Absorbed(a, b)^\times$}{Absorption operator on neighboring branches}

  % ---- Knots, braids, and satellites ----
  \symbolentry{K}{$K \subset S^3$}{A knot}
  \symbolentry{Tpq}{$T_{p,q}$}{The $(p,q)$-torus knot}
  \symbolentry{Trefoil}{$\trefoil$}{The right-handed trefoil knot $T_{2,3}$}
  \symbolentry{Granny}{$\granny$}{The granny knot, connect-sum of two trefoils}
  \symbolentry{Cab}{$\cable{p}{q}(K)$}{The $(p,q)$-cable of $K$ with respect to the Seifert framing}
  \symbolentry{Bn}{$B_n$}{The braid group on $n$ strands}
  \symbolentry{Burau}{$\rho_{n}$}{The reduced Burau representation of $B_n$}

  % ---- Polynomials ----
  \symbolentry{Delta}{$\Alex{K}(t)$}{The Alexander polynomial of $K$}
  \symbolentry{Phi}{$\Cycl{n}(t)$}{The $n$-th cyclotomic polynomial}
  \symbolentry{Phi12}{$\Cycl{12}(t)$}{$t^{4}-t^{2}+1$, governing $\cable{2}{1}(\trefoil)$}

  % ---- Heegaard / knot Floer homology ----
  \symbolentry{HFhat}{$\HF(Y)$}{Hat Heegaard Floer homology of a $3$-manifold $Y$}
  \symbolentry{HFKhat}{$\HFK(K)$}{Hat knot Floer homology of a knot $K$ with $\Ftwo$ coefficients}
  \symbolentry{HFKbi}{$\HFK_{d}(K,s)$}{Bigraded summand at Maslov grading $d$ and Alexander grading $s$}
  \symbolentry{tau}{$\tau(K)$}{The Ozsv\'ath--Szab\'o concordance invariant}
  \symbolentry{Spinc}{$\Spin(Y)$}{The set of $\Spin$ structures on $Y$}

  % ---- Fractional Dehn twist coefficient ----
  \symbolentry{FDTC}{$\FDTC(K)$}{Fractional Dehn twist coefficient of a fibered knot $K$}

  % ---- Miscellaneous ----
  \symbolentry{rank}{$\rank A$}{Rank of an abelian group or $\Ftwo$-vector space $A$}
  \symbolentry{inv}{$\inv{g}$}{The inverse of $g$}
  \symbolentry{ident}{$\id$}{The identity map}
  \symbolentry{pf}{$\pitchfork$}{Transverse intersection}

\end{abbrv}
}
